\par\Needspace{4\baselineskip}
\begin{center}
\SourceLeaf{029}\pdfbookmark[2]{11.}{article-011}11.
\end{center}
The inverse problem, widely known as \textsc{Kepler}\emph{'s problem}, is to find the true anomaly and the radius vector from the mean anomaly. It arises far more often in practice. Astronomers usually express the equation of centre as an infinite series in the sines of the angles $M$, $2M$, $3M$, etc., in which each sine coefficient is itself an infinite series in powers of the eccentricity. Several authors have developed this formula for the equation of centre. We have even less reason to dwell on it here because, in our judgement, it is much less suited to practical use than the indirect method, especially if the eccentricity is anything more than very small. We shall therefore explain the indirect method in somewhat greater detail, in the form we find most convenient.

Equation XII, $E = M + e\sin E$, is transcendental and admits no solution by a finite number of direct operations. It must be solved by successive trials, starting from some approximate value of $E$ and correcting it repeatedly by suitable methods until it satisfies the equation exactly, that is, either to the full precision allowed by the tables of sines or at least to the precision sufficient for the purpose at hand. If these corrections follow a reliable and definite rule, there is scarcely any essential difference between such an indirect method and a solution by series, except that the first value of the unknown is to some extent arbitrary in the indirect method. This may more properly be regarded as an advantage, since a well-chosen value can substantially accelerate the corrections. Let $\epsilon$ be an approximate value of $E$, and let $x$ be the correction still to be added to it (expressed in arcseconds), so that the value $E = \epsilon + x$ satisfies our equation exactly. Compute $e\sin\epsilon$ in arcseconds by logarithms. In doing so, note from the tables the change in $\log\sin\epsilon$ for a change of $1^{\prime\prime}$ in $\epsilon$, and the change in $\log e\sin\epsilon$ for a change of one unit in the number $e\sin\epsilon$. Let these changes, disregarding their signs, be $\lambda$ and $\mu$, respectively. It is hardly necessary to point out that both logarithms are assumed to be given to the same number of decimal places. Now suppose $\epsilon$ is already close enough to the true value of $E$ that the changes in the logarithm of the sine from $\epsilon$ to $\epsilon+x$, and the changes in the logarithm of the number from $e\sin\epsilon$ to $e\sin(\epsilon+x)$, can be treated as uniform. We can then plainly set $e\sin(\epsilon+x) = e\sin\epsilon \pm \frac{\lambda x}{\mu}$, with the \SourceLeaf{030}upper sign applying in the first and fourth quadrants and the lower in the second and third. Hence, since $\epsilon+x = M+e\sin(\epsilon+x)$, we obtain $x = \frac{\mu}{\mu\mp\lambda}(M+e\sin\epsilon-\epsilon)$, and the true value of $E$ is given by $\epsilon+x = M+e\sin\epsilon \pm \frac{\lambda}{\mu\mp\lambda}(M+e\sin\epsilon-\epsilon)$, with the signs determined as described. It is also easy to see that, disregarding the sign, $\mu:\lambda = 1:e\cos\epsilon$, and therefore that $\mu > \lambda$ always. It follows that in the first and last quadrants $M+e\sin\epsilon$ lies between $\epsilon$ and $\epsilon+x$, whereas in the second and third $\epsilon+x$ lies between $\epsilon$ and $M+e\sin\epsilon$. This rule can make it easier to keep track of the signs. If the assumed value $\epsilon$ was still too far from the true value for the foregoing assumption to be sufficiently accurate, this method will certainly yield a much closer value. With that value the same operation must be repeated, once more or several times if necessary. It is readily apparent that if the difference between the initial value $\epsilon$ and the true value is regarded as a first-order quantity, the error in the new value is of the second order. Repeating the operation reduces it to the fourth, eighth, and subsequent orders. Moreover, the smaller the eccentricity, the faster the successive corrections converge.
